

\subsection{Trace Operators and Timoshenko Beams}%\label{sec:trace}

The Saint-Venant solution class admits a six-parameter family of stress states, 
yet the passage from three-dimensional displacement to one-dimensional beam 
kinematics is not unique. Different extraction procedures---here called 
\emph{trace operators}---yield different beam theories: Euler-Bernoulli, 
Timoshenko with Cowper's correction, Timoshenko with Romano's energy-conjugate 
strains, and others. A key observation is that all such theories are 
\emph{exact} within the Saint-Venant framework; they differ not in accuracy 
but in how they partition the displacement field into ``beam'' and ``warping'' 
components.

This distinction has practical consequences. The choice of trace determines:
\begin{enumerate}
    \item the effective boundary conditions implied by a one-dimensional 
    constraint such as $\bm{\theta}(0) = \mathbf{0}$,
    \item the structure of the section constitutive matrix relating beam 
    strains to stress resultants, and
    \item the presence or absence of coupling terms that complicate 
    finite element implementation.
\end{enumerate}
No single trace is universally optimal. The framework developed here makes 
these trade-offs explicit and supports implementation of a general 
cross-section model capable of switching between formulations.

The central mathematical object is the \emph{trace matrix}, which encodes 
how the six Ie\c{s}an parameters map to the coefficients of the beam strain 
field. The \emph{warping index}---the dimension of the kernel of this 
matrix---provides a diagnostic: when it vanishes, the six-strain Cosserat 
model can represent the entire Saint-Venant class; when it is positive, 
additional generalized coordinates are required.

\subsubsection{Ie\c{s}an Parameters and the Resultant Map}

Recall from \S\ref{sec:sv} that the Saint-Venant particular solution 
$\hat{\bm{u}}_{\rm P}$ is parameterized by the loading amplitudes 
$(\bm{a}, \bm{b}_\Section)$ appearing in Problems $\mathrm{P}_{\rm I}$ and 
$\mathrm{P}_{\rm I\!I}$. 
The resultant equations \eqref{eq:2-15} and \eqref{eq:2-18} establish that the 
eight nominal parameters 
$(a_\varepsilon, \bm{a}_\Section, a_\vartheta, b_\varepsilon, \bm{b}_\Section, b_\vartheta)$ 
are subject to two constraints: $b_\varepsilon = -\bm{b}_\Section \cdot \CentroidLocation$ 
and $b_\vartheta = 0$. 
The effective parameter space is therefore six-dimensional.

The \emph{Ie\c{s}an parameter vector} is defined as
\begin{equation}
\IesanVector \triangleq \begin{pmatrix} 
a_\varepsilon \\ \bm{a}_\Section \\ a_\vartheta + c_\vartheta \\ \bm{b}_\Section 
\end{pmatrix} \in \mathbb{R}^{6}
\label{eq:def-iesan-params}
\end{equation}
where $c_\vartheta = c_\vartheta(\bm{b}_\Section)$ is the flexure-induced twist 
determined by \eqref{eq:iesan-2-22}. 
The complete Saint-Venant solution consists of the particular solution 
augmented by an arbitrary rigid body motion \eqref{eq:general-solution}, 
where $\bm{\delta} = (\bm{\delta}_u, \bm{\delta}_\theta) \in \mathbb{R}^6$ 
represents the six rigid body modes. 
The stress field, and hence all resultants, depend only on $\IesanVector$.

The resultants are related to the Ie\c{s}an parameters through a linear map:
\begin{equation}
\begin{pmatrix} N \\ T \\ \ShearRxn \\ \FrameBasis \times \bm{M} \end{pmatrix}
= \IesanMatrix \, \IesanVector
\label{eq:resultant-map}
\end{equation}
where the \emph{resultant matrix} $\IesanMatrix \in \mathbb{R}^{6 \times 6}$ 
encodes the Saint-Venant constitutive relations \eqref{eq:2-15} and \eqref{eq:2-18}:
\begin{equation}
\IesanMatrix = -\begin{pmatrix}
EA & EA\CentroidLocation^\top & 0 & \mathbf{0}^\top \\
0 & \mathbf{0}^\top & G\Jo & G\,\FrameBasis \cdot \int_\Section \bm{r} \times 
\big(\nabla\bm{\WarpShearIesan} + \bm{\Pi}_{(b)}\big) \mathrm{d}A \\
\mathbf{0} & \mathbf{0} & \mathbf{0} & E\IntRGoRG  \\
E\IntRoR \CentroidLocation & E\IntRoR & \mathbf{0} & \mathbf{0}
\end{pmatrix}
\label{eq:def-R-matrix}
\end{equation}


\subsubsection{Trace Operators}

\begin{definition}[Trace operator]
\label{def:sv-trace}
Let $\mathcal{S}_{\mathrm{SV}}$ denote the six-dimensional vector space of 
Saint-Venant displacement fields for a given prismatic bar (modulo global rigid 
motions), parameterized by the Ie\c{s}an vector $\IesanVector \in \mathbb{R}^6$.

A \emph{trace operator} for the Saint-Venant class is a linear map
\[
\mathcal{T} : \mathcal{S}_{\mathrm{SV}} \to C^1(I;\mathbb{R}^6),
\qquad
\hat{\bm{u}} \mapsto 
\begin{pmatrix}
\bm{u} \\[2pt] \bm{\theta}
\end{pmatrix},
\]
where $I \subset \mathbb{R}$ is the axial interval, and the image is interpreted 
as a pair of centerline fields
\[
\bm{u} : I \to \mathbb{R}^3,
\qquad
\bm{\theta} : I \to \mathbb{R}^3.
\]

Given a trace $\mathcal{T}$, the associated \emph{beam strain field} is defined 
by the Cosserat relations \eqref{eq:def-linear-align-strain}
\begin{equation}
\bm{\gamma}(x) 
\;=\;
\bm{u}'(x) + \FrameBasis \times \bm{\theta}(x),
\qquad
\bm{\kappa}(x)
\;=\;
\bm{\theta}'(x),
\label{eq:beam-strains-from-trace}
\end{equation}
and we collect these into a six-dimensional strain vector
\[
\bm{e}(x) 
\;\triangleq\;
\begin{pmatrix}
\bm{\gamma}(x) \\[2pt] \bm{\kappa}(x)
\end{pmatrix} \in \mathbb{R}^6.
\]
\end{definition}


\subsubsection{Affine Structure and Trace Matrices}

\begin{proposition}[Affine-in-$x$ structure of traced strains]
\label{prop:affine-strains}
Let $\mathcal{T}$ be a Saint-Venant linear trace. Then for every 
$\IesanVector \in \mathbb{R}^6$, the beam strain field $\bm{e}(x;\IesanVector)$ 
is at most affine in $x$:
\begin{equation}
\bm{e}(x;\IesanVector)
=
\bm{e}_0(\IesanVector) + x\,\bm{e}_1(\IesanVector),
\label{eq:affine-strains}
\end{equation}
where $\bm{e}_0,\bm{e}_1 : \mathbb{R}^6 \to \mathbb{R}^6$ are linear maps.

Moreover, if $\mathcal{T}$ and $\mathcal{T}_{\mathrm{ref}}$ are two Saint-Venant 
linear traces related by a constant change of strain coordinates
\begin{equation}
\bm{e}(x;\IesanVector)
=
\mathbf{G}\,\bm{e}_{\mathrm{ref}}(x;\IesanVector)
\quad
\text{for all }\IesanVector\in\mathbb{R}^6,\; x\in I,
\label{eq:trace-gauge-transform}
\end{equation}
for some invertible $\mathbf{G} \in \mathbb{R}^{6\times 6}$, then the 
coefficient maps transform as
\[
\bm{e}_0 = \mathbf{G}\,\bm{e}_{0,\mathrm{ref}},
\qquad
\bm{e}_1 = \mathbf{G}\,\bm{e}_{1,\mathrm{ref}}.
\]
\end{proposition}

\begin{proof}
The Saint-Venant particular solution has the structure 
$\hat{\bm{u}}_{\rm P} = \hat{\bm{u}}_{\rm P}^{(0)} + x\,\hat{\bm{u}}_{\rm P}^{(1)} 
+ x^2\,\hat{\bm{u}}_{\rm P}^{(2)}$, where each coefficient depends linearly on 
$\IesanVector$. A Saint-Venant linear trace therefore produces centerline fields 
$\bm{u}(x)$, $\bm{\theta}(x)$ that are at most quadratic in $x$. The strain 
relations \eqref{eq:beam-strains-from-trace} involve one $x$-derivative, 
reducing the polynomial degree by one. Hence $\bm{e}(x)$ is at most affine.

Linearity of $\bm{e}_0$ and $\bm{e}_1$ in $\IesanVector$ follows from linearity 
of the trace. The transformation under $\mathbf{G}$ is immediate from 
\eqref{eq:trace-gauge-transform}.
\end{proof}


\begin{definition}[Trace matrices]
\label{def:trace-matrices}
Fix a Saint-Venant linear trace $\mathcal{T}$. The \emph{trace matrices} 
$\TraceDeDpRoot,\TraceDeDpOne \in \mathbb{R}^{6\times 6}$ are the unique matrices 
such that
\begin{equation}
\bm{e}_0(\IesanVector) = \TraceDeDpRoot\,\IesanVector,
\qquad
\bm{e}_1(\IesanVector) = \TraceDeDpOne\,\IesanVector
\quad
\text{for all }\IesanVector\in\mathbb{R}^6.
\label{eq:def-T0-T1}
\end{equation}

The stacked trace matrix is
\begin{equation}
\mathbf{T}
\;\triangleq\;
\begin{pmatrix}
\TraceDeDpRoot \\[2pt]
\TraceDeDpOne
\end{pmatrix}
\in \mathbb{R}^{12\times 6}.
\label{eq:def-Tstack}
\end{equation}
\end{definition}

\begin{definition}[Warping index]
\label{def:warping-index}
The \emph{warping index} of a trace $\mathcal{T}$ is
\begin{equation}
\WarpIndex
\;\triangleq\;
\dim \ker\!\big(\mathbf{T}\big)
\;=\;
6 - \operatorname{rank}\!\big(\mathbf{T}\big).
\label{eq:def-warp-index}
\end{equation}
\end{definition}

The warping index measures information loss: it counts how many directions in 
the Saint-Venant parameter space produce identically vanishing beam strains.

When $\WarpIndex = 0$, the map $\IesanVector \mapsto (\bm{e}_0, \bm{e}_1)$ is 
injective, and a six-strain Cosserat model can represent the entire 
Saint-Venant class. All warping amplitudes can be recovered as linear 
functions of the beam strains.

When $\WarpIndex > 0$, some Saint-Venant modes are invisible to the beam 
strains. A six-degree-of-freedom element cannot be Saint-Venant exact; one 
must either accept a reduced model or enrich the kinematics with at least 
$\WarpIndex$ additional generalized fields.


\subsubsection{Boundary Conditions and Their Physical Interpretation}

A one-dimensional boundary condition of the form
\begin{equation}
\bm{u}(x_0) = \bm{u}_0, \qquad \bm{\theta}(x_0) = \bm{\theta}_0
\label{eq:1d-bc}
\end{equation}
corresponds to the three-dimensional constraint
\begin{equation}
\hat{\bm{u}}(x_0, \bm{r}) = \bm{u}_0 + \bm{\theta}_0 \times \bm{r} 
+ \Phi(\bm{r})\,\bm{\alpha}(x_0),
\label{eq:3d-bc}
\end{equation}
where $\bm{\alpha}(x_0)$ is the warping amplitude determined by compatibility 
with the Saint-Venant solution. The three-dimensional field at $x_0$ is 
constrained to have prescribed section-rigid components $(\bm{u}_0, \bm{\theta}_0)$ 
\emph{in the sense defined by the trace} $\mathcal{T}$, while the warping 
component $\Phi\,\bm{\alpha}(x_0)$ remains dictated by the global solution.

This has a striking consequence: a ``clamped'' condition 
$\bm{u}(0) = \bm{\theta}(0) = \mathbf{0}$ implies different physical constraints 
depending on the trace.

\begin{itemize}
\item Under a \emph{pointwise trace} at the centroid, the constraint 
$\hat{\bm{u}}(0, \mathbf{0}) = \mathbf{0}$ with 
$(\nabla \times \hat{\bm{u}})|_{(0,\mathbf{0})} = \mathbf{0}$ is enforced: 
the centroid is fixed and the infinitesimal rotation there vanishes.

\item Under a \emph{section-average trace}, the constraints 
$\SectionAverage{\hat{\bm{u}}(0,\cdot)} = \mathbf{0}$ and 
$\RotationAverage{\hat{\bm{u}}(0,\cdot)} = \mathbf{0}$ are enforced: 
the mean translation and mean rotation of the section vanish.

\item Under an \emph{energy-conjugate trace}, the work-conjugate displacement 
and rotation measures vanish, which need not coincide with either pointwise 
or averaged quantities.
\end{itemize}

Crucially, none of these conditions eliminates warping at the boundary. The 
warping amplitude $\bm{\alpha}(0)$ is determined by the global solution and 
will generally be nonzero for loadings involving flexure or torsion. The 
warping displacement performs work against the boundary tractions, giving 
rise to an effective compliance that can be misinterpreted as spring-like 
behavior at the support.

This trace-dependence suggests that the choice of trace should be informed 
by the physical nature of the support. A rigid pin through the centroid is 
better modeled by a pointwise trace; a rigid diaphragm bonded to the entire 
section is better modeled by a section-average trace. However, boundary 
fidelity is only one consideration among several---constitutive symmetry, 
warping index, and compatibility with existing finite element formulations 
may also influence the choice.


\subsubsection{Discussion: Why Torsion Hides the Ambiguity}

The asymmetry between torsion and shear-flexure is striking. Textbooks present 
the torsional constant $\Jo$ as uniquely defined, while shear correction 
factors spawn endless debate. Yet the underlying ambiguity is the same---it 
simply hides better in torsion for three reasons.

First, the torsional warping function is axial and independent of $x$. 
Trace differences that shift the rotation field by an $x$-dependent amount 
do not propagate to the twist strain $\kappa_\vartheta = \theta'_\vartheta$, 
because the shift is annihilated by differentiation.

Second, there is no Poisson coupling in torsion. The energy-conjugacy 
calculation is clean: the integral 
$\int \nabla\WarpTwistIesan \cdot \bm{\tau} \, \mathrm{d}A$ vanishes by the 
traction-free boundary condition, so the virtual work reduces to 
$\delta a_\vartheta \cdot T$ with no residual warping contribution. The 
``natural'' twist strain and the ``energy-conjugate'' twist strain coincide.

Shear-flexure enjoys no such simplification. The integral 
$\int \nabla\bm{\WarpShearIesan} \cdot \bm{\tau} \, \mathrm{d}A$ vanishes by 
the same boundary argument, but 
$\int \bm{\Pi}_{(b)}^\top \bm{\tau} \, \mathrm{d}A$ does not. This Poisson-shear 
coupling forces the energy-conjugate shear strain to differ from any 
kinematically-defined average.

Third, the classical treatment of torsion implicitly works in coordinates 
centered at the shear center---the unique point about which torsion and 
flexure decouple. This is itself a choice of trace, one that happens to 
diagonalize certain couplings. The convention is so universal that it feels 
like physics rather than convention.

The torsion problem thus appears unambiguous because (i) warping is 
$x$-independent, (ii) Poisson effects are absent, and (iii) the shear-center 
convention is tacitly assumed. Shear-flexure strips away all three accidents, 
exposing the trace ambiguity in full generality.